\chapter{Elements of relativistic notations}



A 4-vector is a 1D vector with four elements, 
\begin{align} 
a = (a^0, a1, a^2, a^3) \equiv(a^0, {\bm a})
\end{align}
The first element is called the temporal component; the last three elements  form the spatial components.

The space time vector is 
\begin{align} 
x=(ct, x^1, x^2, x^3)
\end{align}
i.e. its temporal component is the time (multiplied by the speed of light $c$) and the spatial part is a normal position vector.

One can associate 4-vectors not only to the space-time, but also to the linear momentum, and other physical quantities. We can also define 4-vectors which ar not related to a physical quantity, but are  notations for quantities appearing in the calculations.


The scalar product between two 4-vectors is defined :
\begin{align} 
a\cdot b=a^0b^0 -a^1 b^1 -a^2 b^2 -a^3 b^3\equiv a^0b^0-{\bm a}\cdot{\bm b}
\end{align}
Note that the scalar product is not positively defined.

We define the metric tensor
\begin{align}
g_{\mu\nu}\equiv g^{\mu\nu}=(+,-,-,-)
\end{align}
and
\begin{align}
A_\mu=g_{\mu\nu}A^\nu,\qquad A^\mu=g^{\mu\nu}A_\nu,\qquad
\end{align}
